% s-218-p218-m002; original PDF 218, 219
\Aphorism{s-218-p218-m002}{311}
Ich komme hinter eine Bedenklichkeit, die allem Schreiben und Schriftstellern anhaftet: gesetzt, man habe gut, suggestiv wirksam geschrieben, sodass man dem Leser ein gewisses Erleben zur Nachempfindung aufzwingt — gesetzt ferner, man sei zu Zeiten selbst dieser Leser: ist man auch sicher, dass dann noch der Leser den Schreiber aushält? gutheisst? willkommen heisst? Und wenn nicht, wie wehrt man sich gegen sich selbst, der gute Leser gegen den guten Schreiber, — wie schützt man sich vor Infection durch ein früheres, überwundenes Ich, das in aufgezeichneten Seelennöthen unheimlich weiter lebt? Indem man schreibt, wird man sich los, vom Gegenwärtigen an eine Zukunft; aber wer steht dafür, dass die Zukunft die Last williger trägt? Man hat sein psychisches Gift in Tinte umgesetzt; wie aber, wenn die Tinte sich in Gift zurückverwandelt, in abgelagertes, böseres, giftigeres Gift? — Man sollte nur über das ganz Abgeschlossene schreiben, aber — man täusche sich nicht! das Schreiben selber ist noch kein Abschluss. Niemand ist etwas Ernstliches durch Dichten losgeworden, höchstens dichtete er aus Freude, es losgeworden zu sein — und freute sich wahrscheinlich umsonst.
\par

